
% Abstract File
% Do not remove centering environment below
\begin{center}

\end{center}

\begin{center}
{{\bf\fontsize{14pt}{14.5pt}\selectfont \uppercase{ABSTRACT}}}
\end{center}

\doublespacing
\addcontentsline{toc}{chapter}{Abstract}



\begin{center}
	\begin{doublespace}
{\fontsize{14pt}{14.5pt}\selectfont {Title of Dissertation/Thesis}}\\
		{\fontsize{14pt}{14.5pt}\selectfont {Melissa Maria Rios Zertuche}}\\
%		{\fontsize{12pt}{12.5pt}\selectfont {Ph.D. Year}}\\
	\end{doublespace}
\end{center}


% Phages are ubiquitous in nature, and bacteria have evolved a pletora/many defense mechanisms to clear or prevent the spread of viral infections across the population. Among these phage defense mechanisms are retrons. Since their serendipitous discovery, by oversbing spontanous production of ssDNA molecules in Myxoccocuss Xanthus bacteria in 1984, their role in bacteria and mechanisms remained a mistery/not elucidated for over 30 years until recently in 2020, when they were linked to abortive infection pathways triggered by RT, ncRNA with msr and msd (being this one retrotranscribed) and effector proteins joining the vast recently added to the characterized catalog artillery against virus. Since then, characterization and our understating of retron systems slowly began and unlocked interesting/outstanding biotechnological applications including: in situ DNA cargo delivery for genomic engineering by coupling it to CRISPR cas systems, biosensing using ssDNA as molecular barcodes and recording of molecular events. These applications are unlocked by modifying/engineering the ncRNA msd cargo which is co-transcribed with RT to be transcribed depending on the apllicaiton of interest but almost no work on the RT part. To unlock/accelerate the field there are several limitations that has to be addrssed: (i) Data scarcity, (ii) understanding the binding mechanisms behind RT-ncRNA reaction and specificity (iii) determine designibng rules for both RT and ncRNA. Such limitations can be tackled by naturally, expanding the available data for retron systems on large scale data, strenghten or employing more efficient or sofisticated methods and [idea?]. 

% In this thesis we systematically assess the aformentioned limitations by deplyoing a discovery and annotation pipeline on large scale genomic and metogenic data of aproximately 3.9 M entries across different public databases, challenge the current methods and classification of retrons,  address some of the limitation of current traditional methods as data circularity and enhancning ways to detect retron related ncRNA sequences from DNA. By analysing []



Phages are ubiquitous biological entities that exert strong selective pressure on bacterial populations, driving the evolution of diverse mechanisms for the detection and elimination of viral infections. Among these defense systems are retrons, a class of bacterial genetic elements composed of a reverse transcriptase (RT), a structured non-coding RNA (ncRNA), and, in many systems, an associated effector protein. Retrons were first discovered in 1984 through the observation of spontaneous production of single-stranded DNA (msDNA) in \textit{Myxococcus xanthus}; however, their biological function remained largely enigmatic for more than three decades. Their role in bacterial defense was established only recently, when retrons were shown to participate in abortive infection pathways in which the RT uses the retron ncRNA as a template to produce msDNA, ultimately triggering an effector-mediated response to phage infection. This discovery established retrons as an important component of the expanding repertoire of bacterial antiviral defense systems and stimulated increasing interest in their molecular mechanisms and biotechnological potential.

Beyond their biological significance, retrons have emerged as promising platforms for biotechnology. Their ability to generate defined single-stranded DNA molecules \textit{in vivo} has been exploited for applications including \textit{in situ} DNA cargo delivery and genome engineering, biosensing through molecular barcodes, and recording of cellular and molecular events. Most existing engineering efforts, however, have focused primarily on modifying the msd region of the retron ncRNA to alter the sequence of the produced DNA cargo, while comparatively little is known about the RT component itself, particularly the molecular determinants governing RT--ncRNA recognition, compatibility, and specificity. Progress toward predictable retron engineering is therefore limited by three major challenges: (i) the scarcity of experimentally characterized retron systems, (ii) the limited understanding of the sequence, structural, and evolutionary determinants underlying RT--ncRNA recognition, and (iii) the absence of well-defined design principles for engineering both components of the system.

In this thesis, we address these limitations through a systematic computational characterization of retron systems at large scale. We develop and deploy a discovery and annotation pipeline across approximately 4.0 million genomic and metagenomic entries from multiple public databases, enabling the identification and characterization of retron-associated reverse transcriptases and their cognate ncRNAs beyond the experimentally characterized  and previously computationally predicted repertoire. This large-scale analysis expands the available retron sequence space while providing a framework for evaluating the diversity, architecture, and evolutionary relationships of these systems. We further reassess current approaches for retron identification and classification, with particular emphasis on limitations arising from incomplete or circular datasets and on improving the detection and delineation of retron-associated ncRNAs directly from genomic sequence.

Building on this expanded dataset, we investigate the relationship between RT sequence, ncRNA sequence, and ncRNA structure to determine whether information about RT--ncRNA compatibility and recognition can be recovered from their respective molecular representations. By integrating sequence embeddings, structural information, and phylogenetic analyses, we distinguish potential signals of molecular recognition from correlations arising from shared evolutionary history. These analyses provide new insights into the determinants of retron RT--ncRNA association and establish a computational framework for evaluating predictive and design principles for retron systems. Together, this work expands the known diversity of retrons, strengthens their computational discovery and annotation, and provides a foundation for understanding and ultimately engineering RT--ncRNA interactions for future biological and biotechnological applications.


[what to mention about the DL part? mention overview of the results and our labs goal and solution]

